\documentclass[11pt]{article}
\usepackage[margin=1in]{geometry}
\usepackage{amsmath,amssymb,amsthm}
\usepackage{graphicx}
\usepackage{microtype}
\usepackage{xcolor}
\usepackage[hidelinks]{hyperref}

\newtheorem{theorem}{Theorem}
\newtheorem{lemma}[theorem]{Lemma}
\newcommand{\Hh}{\mathbb H}
\newcommand{\Ecal}{\mathcal E}
\newcommand{\R}{\mathbb R}

\title{The Nonnegative Heisenberg Kernel Threshold Is Exactly Four}
\author{Candidate full counterexample to arXiv:1911.03223}
\date{11 August 2026}

\begin{document}
\maketitle

\noindent\fbox{\parbox{0.94\textwidth}{
\textbf{Status: candidate full counterexample, likely valid.}
For every $0<\alpha<4$ there is a compact regular curve in the first
Heisenberg group on which the singular integral associated with $K_\alpha$
is not uniformly $L^2$ bounded.  Hence the positive theorem for
$\alpha\geq4$ is sharp, and in particular it does not persist for
$\alpha\geq2$.  Independent expert review is still required.}}

\section{Source question and current boundary}

F\"assler and Orponen \cite{FO} consider the nonnegative homogeneous kernels
\[
 K_\alpha(x,y,t)
 =\frac{|t|^{\alpha/2}}{\|(x,y,t)\|^{\alpha+1}},\qquad \alpha\geq1,
\]
on the first Heisenberg group.  Their Theorem 1.7 proves boundedness on every
regular curve for $\alpha\geq4$, and the paragraph immediately afterward asks
whether the result persists for $\alpha\geq2$.

\begin{figure}[ht]
\centering
\includegraphics[width=\textwidth]{figures/open_problem_crop.png}
\caption{Theorem 1.7 and the question in F\"assler--Orponen
\cite{FO}, PDF page 4.}
\end{figure}

The May 2026 paper of Chousionis, Li, and Zhang \cite{CLZ} gives
counterexamples for $0<\alpha<2$ and explicitly says that extension to
$[2,4)$ may be possible.  The construction below covers all $0<\alpha<4$.

We use the group law
\[
 (x,y,t)(x',y',t')=(x+x',y+y',t+t'+\tfrac12(xy'-yx'))
\]
and the Kor\'anyi norm
\[
 \|(x,y,t)\|=\big((x^2+y^2)^2+16t^2\big)^{1/4}.
\]
Changing the harmless normalization constant does not affect the argument.

\section{Idea of the proof}

The vertical coordinate of the horizontal lift of a planar graph is exactly
the trapezoidal-rule error for its slope.  We arrange that slope as the
primitive of a bounded stopped dyadic martingale.  Its increments
$\varepsilon_j$ can be square summable---which keeps the martingale bounded
on most paths and the lift regular---while
$\sum_j\varepsilon_j^{\alpha/2}$ diverges whenever $\alpha<4$.

On a fixed positive proportion of pairs crossing the two children of every
active dyadic interval, the current martingale increment dominates all later
increments.  These pair rectangles are disjoint across scales.  Since
$K_\alpha$ is nonnegative, their contributions add rather than cancel, and
the $K_\alpha$ energy is infinite.  Uniform $L^2$ bounds for the truncations
would bound that energy, giving the contradiction.

\section{A stopped-martingale energy lemma}

For $g\in L^\infty[0,1]$ and $y<x$, write
\[
 \Ecal_g(x,y):=\int_y^x\left(\frac{x+y}{2}-s\right)g(s)\,ds.
\]

\begin{lemma}\label{lem:mart}
Let $0<q<2$.  There are a bounded function $g$ on $[0,1]$, decreasing
numbers $\varepsilon_j>0$, dyadic levels $n_j=8j$, and sets $A_j$ which are
unions of dyadic intervals of level $n_j-1$, such that
\[
 |A_j|\geq\frac34,\qquad
 \sum_j\varepsilon_j^2<\infty,\qquad
 \sum_j\varepsilon_j^q=\infty.
\]
Moreover, if $I=[a,a+\ell]\subset A_j$ is one of those level-$(n_j-1)$
intervals, then
\[
 |\Ecal_g(x,y)|\geq c_*\varepsilon_j\ell^2
 \tag{3.1}
\]
whenever
\[
 y\in[a,a+\ell/8],\qquad
 x\in[a+7\ell/8,a+\ell],
\]
where $c_*>0$ is absolute.  All the resulting pair rectangles, over every
$I$ and $j$, are mutually disjoint.
\end{lemma}

\begin{proof}
Choose
\[
 \frac12<p<\frac1q,\qquad \varepsilon_j=cj^{-p},
\]
where $c>0$ is small enough that
$4\sum_j\varepsilon_j^2\leq1/4$.  The two asserted series then have the
required behavior.

Let $\mathcal F_n$ be the dyadic filtration on $[0,1)$, and let $r_n$ be the
$n$th Rademacher function: it is $+1$ and $-1$ on the two children of every
level-$(n-1)$ interval.  Starting with $M_0=0$ and $A_1=[0,1)$, define
\[
 d_j=\varepsilon_j\mathbf 1_{A_j}r_{n_j},\qquad
 M_j=M_{j-1}+d_j,\qquad
 A_{j+1}=A_j\cap\{|M_j|<1/2\}.
\]
This is a stopped martingale.  Doob's inequality gives
\[
 |A_j^c|
 \leq \mathbb P\!\left(\max_{k<j}|M_k|\geq\frac12\right)
 \leq4\mathbb E|M_{j-1}|^2
 \leq4\sum_k\varepsilon_k^2\leq\frac14.
\]
The martingale converges almost everywhere and in $L^2$ to a function $g$.
It is bounded: after the first crossing of $1/2$ the process stops, and the
overshoot is at most $\varepsilon_1$; a path which never crosses stays below
$1/2$.

Fix an active parent $I$ at stage $j$.  Earlier details are constant on $I$,
and $\Ecal$ annihilates constants.  After rescaling $I$ to $[0,1]$, the
contribution of the current Haar detail to $\Ecal_g(x,y)$ has absolute value
\[
 \varepsilon_j\ell^2\frac{(2X-1)(1-2Y)}4
 \geq\frac9{64}\varepsilon_j\ell^2,
 \tag{3.2}
\]
where $Y\in[0,1/8]$ and $X\in[7/8,1]$.

It remains to control descendants.  A later detail $d_k$ has mean zero on
each of its dyadic parent cells $J$, whose length is
$2^{-8(k-j)}\ell$.  On a cell wholly contained in $[y,x]$, subtract a
constant from the affine weight $(x+y)/2-s$; its oscillation on $J$ is at
most $|J|$.  The full cells therefore contribute in total at most
$\varepsilon_k\ell|J|$.  The two possible boundary cells obey the same
bound up to a factor two.  Consequently
\[
 \left|\int_y^x\left(\frac{x+y}{2}-s\right)d_k(s)\,ds\right|
 \leq3\varepsilon_k\ell^2 2^{-8(k-j)}.
\]
Since $\varepsilon_k\leq\varepsilon_j$, the complete descendant tail is at
most
\[
 \frac3{2^8-1}\varepsilon_j\ell^2
 <\frac12\cdot\frac9{64}\varepsilon_j\ell^2.
\]
This proves (3.1), for example with
$c_*=9/64-3/255$.

Finally, a selected pair lies in opposite immediate children of $I$.
Thus $I$ is its unique smallest dyadic common ancestor, which proves
disjointness over all parents and stages.
\end{proof}

\section{The regular curve and divergent energy}

\begin{theorem}\label{thm:main}
For every $0<\alpha<4$, there is a compact regular curve
$\Gamma_\alpha\subset\Hh$ such that the truncations of the singular integral
associated with $K_\alpha$ are not uniformly bounded on
$L^2(\mathcal H^1|_{\Gamma_\alpha})$.
\end{theorem}

\begin{proof}
Set $q=\alpha/2<2$ and take $g$ from Lemma \ref{lem:mart}.  Extend $g$ by
zero outside $[0,1]$, and define
\[
 f(x)=\int_0^x g(u)\,du,\qquad
 F(x)=\int_0^x f(u)\,du.
\]
Since $g$ is bounded, $f$ is Lipschitz.  Consider the horizontal graph map
\[
 \Phi(x)=\left(x,f(x),-F(x)+\frac{x f(x)}2\right).
 \tag{4.1}
\]
For $h=x-y$, direct use of the group law gives
\[
 \Phi(y)^{-1}\Phi(x)
 =\big(h,f(x)-f(y),-\Ecal_g(x,y)\big),
 \tag{4.2}
\]
because integration by parts yields
\[
 F(x)-F(y)-\frac h2(f(x)+f(y))
 =\int_y^x\left(\frac{x+y}{2}-s\right)g(s)\,ds.
\]
The same identity and boundedness of $g$ show
\[
 |f(x)-f(y)|\lesssim|h|,\qquad
 |\Ecal_g(x,y)|\lesssim |h|^2.
\]
Hence
\[
 |x-y|\leq d(\Phi(x),\Phi(y))\lesssim|x-y|.
\]
Thus $\Phi$ is bi-Lipschitz, and
$\Gamma_\alpha:=\Phi([0,1])$ is a compact regular curve.  Its Hausdorff
measure is comparable, under $\Phi$, to Lebesgue measure.

Equation (4.2) and the upper metric estimate imply
\[
 K_\alpha(\Phi(y)^{-1}\Phi(x))
 \gtrsim
 \frac{|\Ecal_g(x,y)|^{\alpha/2}}{|x-y|^{\alpha+1}}.
 \tag{4.3}
\]
For a pair in the rectangle associated with an active interval $I$ of
length $\ell$ at stage $j$, Lemma \ref{lem:mart} and $|x-y|\leq\ell$ give
\[
 K_\alpha(\Phi(y)^{-1}\Phi(x))
 \gtrsim \frac{\varepsilon_j^{\alpha/2}}{\ell}.
\]
The pair rectangle has planar measure $\ell^2/64$, so its contribution to
the positive double integral is at least
$c\varepsilon_j^{\alpha/2}\ell$.  The rectangles are disjoint, and the
total length of the active parents is $|A_j|\geq3/4$.  Consequently
\[
 \iint_{[0,1]^2}
 K_\alpha(\Phi(y)^{-1}\Phi(x))\,dx\,dy
 \gtrsim \sum_j\varepsilon_j^{\alpha/2}|A_j|=\infty.
 \tag{4.4}
\]

To phrase the last step directly in terms of truncations, retain only the
first $N$ generations of pair rectangles and choose the truncation radius
smaller than their least metric separation.  If the truncations were
uniformly $L^2$ bounded, positivity and Cauchy--Schwarz, applied to the
indicator of $\Gamma_\alpha$, would bound every such partial double integral
by a constant times $\mathcal H^1(\Gamma_\alpha)$.  The lower bounds in
(4.4) tend to infinity.  This contradiction proves the theorem.
\end{proof}

\section{Consequences, limitations, and verification}

Theorem \ref{thm:main} answers the source's $\alpha\geq2$ question
negatively, including every previously unsettled value $2\leq\alpha<4$.
Combined with F\"assler--Orponen's positive theorem for $\alpha\geq4$, it
shows that four is the exact threshold.

This does \emph{not} by itself settle the separate smooth-kernel Question 1
shown at the bottom of Figure 1.  The nonsmooth behavior of $K_\alpha$ across
the horizontal plane is essential to membership in this particular family.

The finite checker supplied with the packet verifies the exact Haar formula,
the exponent window, the variance budget, and finite trapezoidal lower bounds.
It is not used in the proof.  A bounded novelty search used the exact formula,
question text, source title and authors, and the terms ``nonnegative kernel'',
``regular curve'', and ``$\alpha=4$''.  It found \cite{CLZ}, which handles
only $0<\alpha<2$ and explicitly flags $[2,4)$ as a possible extension; no
primary source stating the result above was found.

Human review should concentrate on the descendant estimate in Lemma
\ref{lem:mart} and on the Hausdorff-measure transfer in (4.4).

\begin{thebibliography}{9}
\bibitem{FO}
K. F\"assler and T. Orponen,
\emph{Singular integrals on regular curves in the Heisenberg group},
arXiv:1911.03223; J. Math. Pures Appl. 153 (2021), 30--154.

\bibitem{CLZ}
V. Chousionis, S. Li, and L. Zhang,
\emph{On singular integrals with non-negative kernels in the Heisenberg group},
arXiv:2605.17680 (2026).
\end{thebibliography}

\end{document}
